\documentclass[10pt,a4paper]{article}

% --- Packages ---
\usepackage[utf8]{inputenc}
\usepackage[T1]{fontenc}
\usepackage{lmodern}
\usepackage{amsmath,amssymb,amsthm}
\usepackage{mathtools}
\usepackage{geometry}
\usepackage{xcolor}
\usepackage{tcolorbox}
\tcbuselibrary{skins,breakable}
\usepackage{tikz}
\usetikzlibrary{calc,positioning,shapes.geometric,arrows.meta}
\usepackage{fancyhdr}
\usepackage{enumitem}
\usepackage{booktabs}
\usepackage{marginnote}
\usepackage{microtype}
\usepackage{hyperref}

% --- Color Palette ---
\definecolor{primary}{HTML}{0F172A}    % Deep slate blue
\definecolor{secondary}{HTML}{0284C7}  % Ocean blue
\definecolor{accent}{HTML}{0D9488}     % Teal accent
\definecolor{bglight}{HTML}{F8FAFC}    % Light gray fill
\definecolor{bordercol}{HTML}{CBD5E1} % Slate border
\definecolor{textdark}{HTML}{334155}  % Dark charcoal body text
\definecolor{sidenotecol}{HTML}{475569} % Muted slate for side notes

% --- Page Setup ---
\geometry{
  a4paper,
  top=14mm,
  bottom=14mm,
  left=14mm,
  right=56mm, % Generous right margin for side notes
  marginparwidth=46mm,
  marginparsep=4mm,
  headheight=18pt
}

\pagestyle{fancy}
\fancyhf{}
\renewcommand{\headrulewidth}{0.4pt}
\renewcommand{\headrule}{{\color{bordercol}\hrule width\headwidth height 0.4pt}}
\fancyhead[L]{\small\bfseries\color{primary}MATH 204: Calculus \& Linear Algebra}
\fancyhead[R]{\small\color{sidenotecol}Worksheet 04 --- Fall 2026}
\fancyfoot[C]{\small\color{sidenotecol}Page \thepage\ of 3}

% --- Custom Fonts \& Settings ---
\renewcommand{\familydefault}{\sfdefault}
\usepackage{sansmath}
\sansmath

\newcommand{\sidenote}[1]{%
  \marginnote{\raggedright
    \begin{tcolorbox}[colframe=secondary!30,colback=bglight,boxrule=0.5pt,arc=2pt,left=4pt,right=4pt,top=2pt,bottom=2pt]
      \fontsize{8.5}{11}\selectfont\color{sidenotecol}\textbf{\color{secondary}Note:} #1
    \end{tcolorbox}%
  }%
}

\newcommand{\conceptbox}[2]{%
  \marginnote{\raggedright
    \begin{tcolorbox}[colframe=accent!40,colback=accent!5,boxrule=0.6pt,arc=3pt,left=4pt,right=4pt,top=2pt,bottom=2pt]
      \fontsize{8.5}{11}\selectfont\color{textdark}\textbf{\color{accent}#1}\\ #2
    \end{tcolorbox}%
  }%
}

\newtcolorbox{problembox}[1]{
  enhanced,
  colframe=primary,
  colback=bglight,
  coltitle=white,
  fonttitle=\bfseries\sffamily,
  title={#1},
  boxrule=0.8pt,
  arc=3pt,
  top=2.5pt,bottom=2.5pt,left=6pt,right=6pt,
  shadow={1pt}{-1pt}{0pt}{opacity=0.08,black}
}

\newtcolorbox{solutionbox}{
  enhanced,
  colframe=secondary!40,
  colback=white,
  boxrule=0.5pt,
  arc=2pt,
  top=2.5pt,bottom=2.5pt,left=6pt,right=6pt,
  borderline west={3pt}{0pt}{secondary}
}

\newtcolorbox{keyconcept}{
  enhanced,
  colframe=accent,
  colback=accent!5,
  boxrule=0.6pt,
  arc=4pt,
  title={\textbf{Key Result / Theorem Summary}},
  coltitle=white,
  fonttitle=\small\bfseries,
  top=2.5pt,bottom=2.5pt,left=6pt,right=6pt
}

\setlength{\parindent}{0pt}
\setlength{\parskip}{1.5pt}

\begin{document}

% --- Header Banner ---
\begin{tcolorbox}[
  enhanced,
  colback=primary,
  colframe=primary,
  arc=4pt,
  top=5pt,bottom=5pt,left=10pt,right=10pt
]
  \color{white}
  \begin{minipage}[c]{0.65\textwidth}
    {\Large \bfseries Apex Institute of Technology}\\[1pt]
    {\small Department of Mathematics --- MATH 204}\\[2pt]
    {\LARGE \bfseries Multivariable Calculus \& Linear Algebra}
  \end{minipage}%
  \hfill
  \begin{minipage}[c]{0.32\textwidth}
    \raggedleft \small
    \textbf{Worksheet 04}\\
    \textbf{Topic:} Spectral Theorem \& Vector Fields\\
    \textbf{Term:} Fall 2026\\
    \textbf{Instructor:} Dr. Elena Vance
  \end{minipage}
\end{tcolorbox}

\vspace{0.5mm}

\begin{tcolorbox}[colback=bglight,colframe=bordercol,boxrule=0.5pt,arc=2pt,top=2pt,bottom=2pt,left=8pt,right=8pt]
  \small \color{textdark}
  \textbf{Student Name:} \underline{\hspace{4.2cm}} \hfill
  \textbf{Student ID:} \underline{\hspace{2.8cm}} \hfill
  \textbf{Date:} \underline{\hspace{2.2cm}}
\end{tcolorbox}

% ==============================================================================
% SECTION 1: LINEAR ALGEBRA
% ==============================================================================
\section*{\color{primary}\large 1. Linear Transformations \& Eigenvalue Problems}

\sidenote{A real symmetric matrix $A = A^T$ has orthogonal eigenvectors for distinct eigenvalues. This guarantees diagonalizability over $\mathbb{R}$.}

\begin{problembox}{Problem 1.1: Symmetric Matrix Diagonalization \& Spectral Decomposition}
Consider the symmetric matrix $A \in \mathbb{R}^{3 \times 3}$ defined by:
\[
A = \begin{bmatrix} 3 & 1 & 1 \\ 1 & 3 & 1 \\ 1 & 1 & 3 \end{bmatrix}
\]
\begin{enumerate}[label=(\alph*),leftmargin=*]
  \item Find the characteristic polynomial $p_A(\lambda) = \det(A - \lambda I)$ and derive all eigenvalues.
  \item Construct an orthonormal basis of eigenvectors for $\mathbb{R}^3$ and matrix $P$ such that $P^T A P = D$.
  \item Express $A^5$ using the spectral decomposition $A = \sum_{i} \lambda_i \mathbf{u}_i \mathbf{u}_i^T$.
\end{enumerate}
\end{problembox}

\conceptbox{Algebraic Multiplicity}{Multiplicity of $\lambda$ as a root of $\det(A - \lambda I) = 0$. Geometric multiplicity satisfies $1 \le \dim(\ker(A - \lambda I)) \le \text{alg. mult.}$}

\begin{solutionbox}
\textbf{Detailed Derivation Step-by-Step:}

\textbf{Step 1: Compute eigenvalues.} Subtract $\lambda$ along the main diagonal:
\[
\det(A - \lambda I) = \det \begin{bmatrix} 3-\lambda & 1 & 1 \\ 1 & 3-\lambda & 1 \\ 1 & 1 & 3-\lambda \end{bmatrix} = (2-\lambda) \det \begin{bmatrix} 1 & 0 & 0 \\ 1 & 4-\lambda & 1 \\ 1 & 2 & 3-\lambda \end{bmatrix}
\]
Expanding the remaining $2 \times 2$ determinant:
\[
(4-\lambda)(3-\lambda) - 2 = \lambda^2 - 7\lambda + 10 = (\lambda - 2)(\lambda - 5) \implies p_A(\lambda) = -(\lambda - 2)^2 (\lambda - 5) = 0
\]
Thus, the eigenvalues are $\lambda_1 = 5$ (mult. 1) and $\lambda_2 = \lambda_3 = 2$ (mult. 2).

\textbf{Step 2: Eigenspace for $\lambda_1 = 5$.}
\[
(A - 5I)\mathbf{v} = \mathbf{0} \implies -2x + y + z = 0, \; x - 2y + z = 0 \implies \mathbf{u}_1 = \frac{1}{\sqrt{3}}\begin{bmatrix} 1 \\ 1 \\ 1 \end{bmatrix}
\]

\textbf{Step 3: Eigenspace for $\lambda_2 = 2$.}
\[
(A - 2I)\mathbf{v} = \mathbf{0} \implies x + y + z = 0 \implies \mathbf{u}_2 = \frac{1}{\sqrt{2}}\begin{bmatrix} 1 \\ -1 \\ 0 \end{bmatrix}, \; \mathbf{u}_3 = \frac{1}{\sqrt{6}}\begin{bmatrix} 1 \\ 1 \\ -2 \end{bmatrix}
\]
\end{solutionbox}

\newpage

% ==============================================================================
% PAGE 2: DERIVATION CONTINUED & CALCULUS
% ==============================================================================

\sidenote{Gram-Schmidt process ensures orthogonality when eigenspaces have dimension $\ge 2$. Here $\mathbf{u}_2 \cdot \mathbf{u}_3 = 0$ holds by construction.}

\begin{solutionbox}
\textbf{Step 4: Orthogonal Matrix $P$ and Spectral Power Calculation.}
\[
P = \begin{bmatrix} \frac{1}{\sqrt{3}} & \frac{1}{\sqrt{2}} & \frac{1}{\sqrt{6}} \\[2pt] \frac{1}{\sqrt{3}} & -\frac{1}{\sqrt{2}} & \frac{1}{\sqrt{6}} \\[2pt] \frac{1}{\sqrt{3}} & 0 & -\frac{2}{\sqrt{6}} \end{bmatrix}, \quad D = \begin{bmatrix} 5 & 0 & 0 \\ 0 & 2 & 0 \\ 0 & 0 & 2 \end{bmatrix}
\]
Using projection operators $P_1 = \mathbf{u}_1\mathbf{u}_1^T = \frac{1}{3}\mathbf{J}$ and $P_2 = I - \frac{1}{3}\mathbf{J}$ (where $\mathbf{J}$ is all ones):
\[
A^5 = 5^5 P_1 + 2^5 P_2 = 3125 \left(\frac{1}{3}\mathbf{J}\right) + 32 \left(I - \frac{1}{3}\mathbf{J}\right) = 32 I + 1031 \mathbf{J} = \begin{bmatrix} 1063 & 1031 & 1031 \\ 1031 & 1063 & 1031 \\ 1031 & 1031 & 1063 \end{bmatrix}
\]
\end{solutionbox}

% ==============================================================================
% SECTION 2: VECTOR CALCULUS & LINE INTEGRALS
% ==============================================================================
\section*{\color{primary}\large 2. Multivariable Calculus \& Vector Field Integration}

\conceptbox{Conservative Fields}{A vector field $\mathbf{F} = \nabla f$ is conservative iff $\text{curl}\,\mathbf{F} = \mathbf{0}$ on a simply connected domain. Then $\int_C \mathbf{F} \cdot d\mathbf{r} = f(B) - f(A)$.}

\begin{problembox}{Problem 2.1: Green's Theorem \& Circulation of Vector Fields}
Let $C$ be the boundary of the region $R$ bounded by $y = x^2$ and $y = 4$, oriented counterclockwise. Evaluate the line integral:
\[
I = \oint_C \left( (y^2 + e^{\sqrt{x}}) \, dx + (3xy + \sin(y^2)) \, dy \right)
\]
\end{problembox}

\sidenote{Green's Theorem transforms a closed line integral $\oint_{\partial R} P\,dx + Q\,dy$ into a double integral $\iint_R \left(\frac{\partial Q}{\partial x} - \frac{\partial P}{\partial y}\right) dA$ over region $R$.}

\begin{solutionbox}
\textbf{Step-by-Step Application of Green's Theorem:}

\textbf{1. Compute component partial derivatives:}
\[
P(x,y) = y^2 + e^{\sqrt{x}} \implies \frac{\partial P}{\partial y} = 2y, \qquad Q(x,y) = 3xy + \sin(y^2) \implies \frac{\partial Q}{\partial x} = 3y
\]
\textbf{2. Formulate double integral over region $R$:}
\[
\frac{\partial Q}{\partial x} - \frac{\partial P}{\partial y} = 3y - 2y = y \implies I = \iint_R y \, dA = \int_{-2}^{2} \int_{x^2}^{4} y \, dy \, dx
\]
\textbf{3. Evaluate inner and outer integrals:}
\[
\int_{x^2}^{4} y \, dy = \left[ \frac{y^2}{2} \right]_{x^2}^{4} = 8 - \frac{x^4}{2} \implies I = 2 \int_{0}^{2} \left( 8 - \frac{x^4}{2} \right) dx = 2 \left[ 8x - \frac{x^5}{10} \right]_{0}^{2} = \frac{128}{5} = 25.6
\]
\end{solutionbox}

\newpage

% ==============================================================================
% PAGE 3: PRACTICE PROBLEMS & APPLIED EXERCISES
% ==============================================================================

\begin{keyconcept}
\textbf{Core Theorems of Multivariable Calculus in $\mathbb{R}^2$ and $\mathbb{R}^3$:}

\vspace{0.5mm}
\centering
\fontsize{8}{10}\selectfont
\begin{tabular}{@{}lll@{}}
  \toprule
  \textbf{Theorem} & \textbf{Domain Relation} & \textbf{Integrands Formula} \\
  \midrule
  Fundamental Theorem & Curve $C$ from $A$ to $B$ & $\int_C \nabla f \cdot d\mathbf{r} = f(B) - f(A)$ \\{}
  Green's Theorem & Boundary $\partial R \to R \subset \mathbb{R}^2$ & $\oint_{\partial R} \mathbf{F} \cdot d\mathbf{r} = \iint_R \left(\frac{\partial Q}{\partial x} - \frac{\partial P}{\partial y}\right) dA$ \\{}
  Divergence Theorem & Boundary $\partial R \to R \subset \mathbb{R}^2$ & $\oint_{\partial R} \mathbf{F} \cdot \mathbf{n}\,ds = \iint_R (\nabla \cdot \mathbf{F})\,dA$ \\
  \bottomrule
\end{tabular}
\end{keyconcept}

\section*{\color{primary}\large 3. Practice Problems \& Conceptual Exercises}

\conceptbox{Submission Guidelines}{Complete Problems 3.1 and 3.2 for peer grading. Show all intermediate steps and state relevant theorems explicitly.}

\begin{problembox}{Problem 3.1: Coupled Linear System of Differential Equations}
Consider the initial value problem for a coupled system of ODEs:
\[
\frac{d\mathbf{x}}{dt} = A\mathbf{x}(t), \quad \mathbf{x}(0) = \begin{bmatrix} 2 \\ 1 \end{bmatrix}, \quad \text{where } A = \begin{bmatrix} 1 & 4 \\ 2 & 3 \end{bmatrix}
\]
\begin{enumerate}[label=(\alph*),leftmargin=*]
  \item Find the eigenvalues and eigenvectors of $A$, and determine the matrix exponential $e^{At}$.
  \item Solve for the exact trajectory $\mathbf{x}(t)$ and analyze stability near the origin.
\end{enumerate}
\end{problembox}

\begin{problembox}{Problem 3.2: Outward Flux Across a Closed Surface (Divergence Theorem)}
Let $\mathbf{F}(x,y,z) = (x^3 + y)\mathbf{i} + (y^3 + z)\mathbf{j} + (z^3 + x)\mathbf{k}$.
Compute the outward flux $\iint_S \mathbf{F} \cdot d\mathbf{S}$ across boundary $S$ of sphere $B: x^2 + y^2 + z^2 \le 1$.
\end{problembox}

\sidenote{Notice $\nabla \cdot \mathbf{F} = 3(x^2 + y^2 + z^2) = 3\rho^2$. Switch to spherical coordinates $(\rho, \phi, \theta)$ to evaluate the triple integral efficiently.}

\begin{solutionbox}
\textbf{Outline Solution for Problem 3.2:}
By Gauss's Divergence Theorem:
\begin{align*}
\Phi &= \iint_S \mathbf{F} \cdot d\mathbf{S} = \iiint_B (\nabla \cdot \mathbf{F}) \, dV = 3 \iiint_B (x^2 + y^2 + z^2) \, dV \\
&= 3 \int_{0}^{2\pi} d\theta \int_{0}^{\pi} \sin\phi \, d\phi \int_{0}^{1} \rho^4 \, d\rho = 3(2\pi)(2)\left( \frac{1}{5} \right) = \frac{12\pi}{5}
\end{align*}
\end{solutionbox}

\vspace{1mm}

\begin{tcolorbox}[colframe=secondary!40,colback=bglight,boxrule=0.5pt,arc=3pt,top=2pt,bottom=2pt,left=8pt,right=8pt]
  \centering \small \color{sidenotecol}
  \textbf{Apex Institute of Technology} --- Department of Mathematics \\
  MATH 204: Multivariable Calculus \& Linear Algebra --- Compiled with \LaTeX
\end{tcolorbox}

\end{document}
